% Pandoc template for IEEE conference paper sections
\IEEEpubidadjcol
\section{Preliminaries}

\subsection{Elliptic Curves and
Pairings}\label{elliptic-curves-and-pairings}

Let \(\mathbb{F}_p\) be a prime field with \(p > 3\). An elliptic curve
in Short Weierstrass form is \(E: y^2 = x^3 + ax + b\) with
\(4a^3 + 27b^2 \neq 0\). The set \(E(\mathbb{F}_p)\) of rational points
forms an abelian group under the chord-and-tangent law
\cite{Silverman2009}.

A \emph{bilinear pairing} is a map
\(e: \mathbb{G}_1 \times \mathbb{G}_2 \to \mathbb{G}_T\) satisfying:

\begin{itemize}
\tightlist
\item
  \textbf{Bilinearity:} \(e(aP, bQ) = e(P, Q)^{ab}\) for all
  \(a, b \in \mathbb{Z}\)
\item
  \textbf{Non-degeneracy:} \(e(P, Q) \neq 1\) for generators
  \(P \in \mathbb{G}_1\), \(Q \in \mathbb{G}_2\)
\item
  \textbf{Computability:} \(e\) can be evaluated in polynomial time via
  Miller's algorithm \cite{Miller1986}
\end{itemize}

Here \(\mathbb{G}_1, \mathbb{G}_2\) are prime-order subgroups of
\(E(\mathbb{F}_p)\) and \(E(\mathbb{F}_{p^k})\) respectively, and
\(\mathbb{G}_T \subset \mathbb{F}_{p^k}^*\) is the target group. The
principal pairing constructions are the Weil pairing (symmetric,
computes two Miller loops) and the Tate pairing (one Miller loop plus a
final exponentiation). The \emph{optimal Ate pairing}
\cite{Vercauteren2010} further reduces the Miller loop length to
\(|T|\) bits (the curve seed), making it the preferred choice for
implementations.

\subsection{Embedding Degree and
Security}\label{embedding-degree-and-security}

A curve is \emph{pairing-friendly} with embedding degree \(k\) if there
exists a prime-order subgroup of size \(r\) such that \(r \mid p^k - 1\)
but \(r \nmid p^i - 1\) for \(0 < i < k\). The embedding degree
determines the pairing target field \(\mathbb{F}_{p^k}\): the MOV attack
\cite{MOV1993} reduces ECDLP in \(\mathbb{G}_1\) to DLP in
\(\mathbb{F}_{p^k}^*\), so \(k\) must be large enough that DLP in
\(\mathbb{F}_{p^k}\) remains intractable.

For a 1024-bit base field with \(k = 18\), the pairing target is
\(\mathbb{F}_{p^{18}}\)---a field of 18,432 bits. This far exceeds the
3,072-bit threshold for 128-bit DLP security and provides substantial
margin against the Tower Number Field Sieve (TNFS)
\cite{BarbulescuDuquesne2019, KimBarbulescu2016}. The choice \(k = 18\)
also admits the 18th cyclotomic polynomial
\(\Phi_{18}(T) = T^6 - T^3 + 1\), which produces 512-bit primes \(r\)
from 86-bit seeds---an efficient ratio for the Cocks-Pinch construction.

\subsection{Complex Multiplication
Method}\label{complex-multiplication-method}

The CM method constructs curves with a prescribed group order. Given a
CM discriminant \(D < 0\), the existence of a curve over
\(\mathbb{F}_p\) with Frobenius trace \(t\) requires the Diophantine
equation:

\[4p = t^2 + |D| \cdot y^2\]

to have an integer solution. The \(j\)-invariant is obtained from the
Hilbert class polynomial \(H_D(x)\). For \(D = -3\), we have
\(H_{-3}(x) = x\), giving \(j = 0\) and the simple curve form
\(y^2 = x^3 + b\).

\subsection{Traditional Cocks-Pinch
Algorithm}\label{traditional-cocks-pinch-algorithm}

The Cocks-Pinch method \cite{CocksPinch2001} constructs
pairing-friendly curves by starting from the subgroup order \(r\) rather
than the field prime \(p\). For embedding degree \(k = 18\):

\begin{enumerate}
\def\labelenumi{\arabic{enumi}.}
\tightlist
\item
  \textbf{Generate \(r\)}: Choose a parameter \(T\) and compute
  \(r = \Phi_{18}(T) = T^6 - T^3 + 1\). Verify that \(r\) is prime.
\item
  \textbf{Compute square root}: Find \(\beta = \sqrt{-3} \pmod{r}\) via
  Tonelli-Shanks.
\item
  \textbf{CM parameters}: For each \(i\) with \(\gcd(i, 18) = 1\),
  compute \(t_0 \equiv T^i + 1 \pmod{r}\) and
  \(y_0 \equiv (t_0 - 2)/\beta \pmod{r}\).
\item
  \textbf{Lift to integers}: Search over \((h_t, h_y) \in [-20, 20]^2\)
  for \(t = t_0 + h_t \cdot r\), \(y = y_0 + h_y \cdot r\) such that
  \(p = (t^2 + 3y^2)/4\) is prime of the target size.
\end{enumerate}

The resulting curve has \(\rho = \log p / \log r \approx 2\), which is
the standard tradeoff of the Cocks-Pinch method compared to parametric
families (\(\rho = 1\) for BN, \(\rho = 1.5\) for BLS12).

\subsection{Two-Adicity and NTT}\label{two-adicity-and-ntt}

The \emph{two-adicity} of a prime \(q\) is \(\nu_2(q-1)\), the largest
power of 2 dividing \(q-1\). A field \(\mathbb{F}_q\) with two-adicity
\(s\) contains a primitive \(2^s\)-th root of unity, enabling NTT of
length up to \(2^s\). Modern ZK proof systems (Groth16
\cite{Groth2016}, PLONK \cite{PLONK2019}) require NTT for polynomial
multiplication over both \(\mathbb{F}_r\) (constraint evaluation) and
\(\mathbb{F}_p\) (pairing computation). BLS12-381 achieves two-adicity
32 in both fields \cite{BLS12_381}.
